El aprendizaje automático o \texttt{Machine Learning} es una rama de la inteligencia artificial cuyo objetivo es el desarrollo de modelos que aprenden patrones a partir de conjuntos de datos. Estos modelos ajustan su comportamiento mediante muestras pertenecientes a estos conjuntos de datos, lo que permite realizar predicciones sobre datos no vistos previamente. En el ámbito de los IDS, esto resulta bastante útil debido a la variabilidad del tráfico de red. Debido a esto, se ha estudiado la aplicación de técnicas de aprendizaje automático en IDS con el objetivo de mejorar la tasa de detección, reducir el número de falsos positivos y aumentar la adaptación del IDS contra nuevos ataques \cite{zamani2015mlids,buczak2016survey}.

En un IDS basado en aprendizaje automático, cada muestra del conjunto de datos representa normalmente un paquete o flujo de estos o una conexión. Cada una de estas muestras se describe mediante un conjunto de características, cada una de las cuales puede aportar información temporal, estadística, de volumen, de protocolo o de comportamiento. A partir de estas, el modelo predice una etiqueta o clase. Dependiendo de cómo se defina la etiqueta, esta puede indicar únicamente si la muestra pertenece a tráfico benigno o malicioso (problema de clasificación binaria) o el tipo específico de ataque (problema de clasificación multiclase). Un problema de clasificación multiclase es más complejo pero aporta más información \cite{wagh2013surveyids}.

Para implementar un IDS basado en aprendizaje automático se siguen varios pasos. En primer lugar, se recopila tráfico que sirva como base experimental. Después, se realiza una fase de preprocesamiento, en la que se limpian los datos y se preparan las características para que puedan ser utilizadas por el modelo. Tras esto, el conjunto de datos se divide en subconjuntos de entrenamiento y prueba. El subconjunto de entrenamiento se utiliza para alimentar el modelo y que pueda aprender patrones en los valores de las características asociadas a cierta etiqueta, mientras que el subconjunto de prueba permite evaluar el entrenamiento y ver si las muestras se clasifican correctamente. Finalmente, el modelo entrenado se evalúa mediante métricas \cite{buczak2016survey}.

Hay varios tipos de aprendizaje automático: aprendizaje supervisado, aprendizaje no supervisado y arprendizaje semisupervisado. El aprendizaje supervisado consiste en entrenar el modelo con muestras etiquetadas, conociendo de antemano la clase a la que pertenecen. Por otro lado, en el aprendizaje no supervisado el modelo no dispone de etiquetas y debe encontrar patrones a partir de la estructura interna de los datos, por lo que su evaluación es más compleja aunque por eso mismo permite detectar ataques desconocidos. Finalmente, en un modelo semisupervisado solo una parte de las muestras dispone de etiquetas. En general, los métodos supervisados obtienen buenos resultados cuando las muestras de la evaluación son similares a las del entrenamiento, mientras que los métodos no supervisados son más útiles cuando se pretende detectar nuevas patrones \cite{zamani2015mlids,kunal2019mlids}.

Entre los modelos supervisados más utilizados en detección de intrusiones se encuentran los clasificadores bayesianos, especialmente Naive Bayes. Estos modelos se basan en el teorema de Bayes y calculan la probabilidad de que una muestra pertenezca a una determinada clase a partir de los valores de sus características. Naive Bayes introduce una simplificación importante al asumir independencia condicional entre las características dada la clase. Esta hipótesis permite construir modelos rápidos y sencillos, pero puede limitar su rendimiento cuando existen fuertes dependencias entre variables, algo frecuente en tráfico de red. Aun así, su bajo coste computacional y su simplicidad hacen que se utilice a menudo como modelo base de comparación en trabajos de IDS \cite{halimaa2019mlbasedids,wagh2013surveyids}.

Otro modelo ampliamente utilizado es \emph{Support Vector Machine} (SVM). Este algoritmo busca construir un hiperplano que separe las muestras de distintas clases maximizando el margen entre ellas. Cuando los datos no son linealmente separables, SVM puede utilizar funciones \emph{kernel}, como funciones polinómicas, sigmoidales o de base radial, para proyectar los datos a un espacio de mayor dimensión en el que la separación entre clases resulte más sencilla. En detección de intrusiones, SVM se ha utilizado tanto para clasificación binaria como para clasificación multiclase, especialmente en escenarios donde las fronteras entre tráfico normal y anómalo no son triviales. Sin embargo, su coste computacional puede crecer cuando el número de muestras es elevado, lo que debe tenerse en cuenta en conjuntos de datos de red de gran tamaño \cite{halimaa2019mlbasedids,zamani2015mlids}.

Los árboles de decisión constituyen otra familia de modelos relevante en IDS. Estos modelos dividen progresivamente el espacio de características mediante reglas de decisión, hasta asignar una clase a cada muestra. Su principal ventaja es que generan estructuras relativamente interpretables, ya que las decisiones pueden representarse como una secuencia de condiciones sobre las variables de entrada. En el contexto de detección de intrusiones, esta interpretabilidad puede ser útil para comprender qué características del tráfico están asociadas a determinados ataques. No obstante, un único árbol de decisión puede ser sensible al ruido y al sobreajuste, especialmente cuando el conjunto de datos contiene muchas características o clases desbalanceadas.

También son frecuentes los métodos de ensamblado, que combinan varios clasificadores para obtener una predicción más robusta que la de un único modelo. A partir de los árboles de decisión surge Random Forest, un método de ensamblado que combina múltiples árboles para mejorar la capacidad de generalización. Cada árbol se entrena introduciendo aleatoriedad, tanto en las muestras utilizadas como en las características consideradas durante las divisiones, y la predicción final se obtiene normalmente mediante voto mayoritario. Esta combinación permite reducir la dependencia respecto a un único árbol y mejorar la robustez del modelo. Además, Random Forest proporciona estimaciones internas del error y permite calcular la importancia de las características, lo que resulta especialmente útil cuando se desea analizar qué variables contribuyen más a la clasificación \cite{breiman2001randomforest}. Por estas razones, Random Forest es un modelo frecuente en problemas de IDS basados en datos tabulares, donde las muestras se representan mediante características extraídas del tráfico de red.

También se han aplicado técnicas basadas en vecinos próximos, como \emph{k-Nearest Neighbors} (kNN). Este modelo clasifica una muestra en función de las clases de sus vecinos más cercanos según una medida de distancia. Su funcionamiento es sencillo y puede adaptarse a problemas de clasificación binaria o multiclase. Sin embargo, en IDS puede presentar limitaciones cuando el conjunto de datos es muy grande, ya que la clasificación de nuevas muestras requiere calcular distancias respecto a numerosos ejemplos de entrenamiento. Además, su rendimiento depende de la escala de las características y de la elección de la medida de distancia, por lo que suele requerir normalización y preprocesamiento adecuado \cite{zamani2015mlids,buczak2016survey}.

En el caso del aprendizaje no supervisado, una técnica habitual es el clustering. Estos métodos agrupan las muestras en función de su similitud, sin utilizar etiquetas durante el entrenamiento. Algoritmos como \emph{k-means} permiten identificar grupos de tráfico con comportamiento parecido, de modo que las muestras alejadas de los patrones normales pueden considerarse potencialmente anómalas. En IDS, el clustering puede ser útil cuando no se dispone de etiquetas fiables o cuando se quiere detectar tráfico que no encaja con el comportamiento habitual. No obstante, la elección del número de grupos, la medida de distancia y la interpretación posterior de los clústeres pueden influir de forma significativa en los resultados \cite{wagh2013surveyids}.

Las redes neuronales artificiales también han sido ampliamente estudiadas en detección de intrusiones. Estos modelos están formados por unidades de procesamiento interconectadas, organizadas en capas, que aprenden relaciones no lineales entre las características de entrada y la clase de salida. Su capacidad para modelar patrones complejos las hace adecuadas para problemas donde la frontera entre tráfico benigno y malicioso no es sencilla. A partir de estas redes surgen los modelos de aprendizaje profundo, o \emph{Deep Learning}, que emplean arquitecturas con múltiples capas para aprender representaciones más complejas de los datos. En IDS se han utilizado redes neuronales profundas, redes convolucionales, redes recurrentes y autoencoders. Las redes convolucionales pueden emplearse para detectar patrones locales en representaciones del tráfico, las redes recurrentes resultan útiles cuando se desea modelar dependencias temporales o secuenciales, y los autoencoders se han aplicado especialmente en detección de anomalías, aprendiendo una representación del tráfico normal y detectando como sospechosas aquellas muestras con alto error de reconstrucción. Estas técnicas pueden alcanzar buenos resultados, especialmente cuando se dispone de grandes volúmenes de datos, pero suelen requerir mayor coste computacional, más tiempo de entrenamiento y presentan menor interpretabilidad que muchos modelos clásicos \cite{wagh2013surveyids,buczak2016survey}.

Un aspecto común a todos estos modelos es la importancia de las características empleadas para representar el tráfico. Los conjuntos de datos de red suelen contener un número elevado de variables, algunas de las cuales pueden ser redundantes, irrelevantes o ruidosas. Esto puede aumentar el tiempo de entrenamiento, dificultar la generalización y empeorar el rendimiento del clasificador. Por ello, muchos trabajos incorporan técnicas de selección de características, reducción de dimensionalidad o ingeniería de características como parte del preprocesamiento. Estas técnicas buscan conservar la información más relevante para la detección de ataques, reduciendo al mismo tiempo la complejidad del modelo \cite{kunal2019mlids,buczak2016survey}.
